% GENERATED FILE. Edit canonical lessons, then run npm run generate:books.
% canonical-source-hash: fnv1a64:f5cc373aff07e2e3
% canonical-lessons: TE-C355-varsikotsavam, TE-C355-pelliroju, TE-C355-sastipurti, TE-C355-jayanti, TE-C355-sanmanam

\chapter{Anniversaries}
\label{ch:te-anniversaries}
\hlchaptermodality{355}

\begin{chapteropening}
\textbf{By the end of this chapter:} \emph{I can say an anniversary, a wedding anniversary, a sixtieth birthday, a birth anniversary and a felicitation.}

Complete the last lesson of chapter 355: I can say an anniversary, a wedding anniversary, a sixtieth birthday, a birth anniversary and a felicitation.
\end{chapteropening}

\section[\texorpdfstring{\te{వార్షికోత్సవం} (vārṣikōtsavaṁ)}{వార్షికోత్సవం (vārṣikōtsavaṁ)}]{\te{వార్షికోత్సవం} (vārṣikōtsavaṁ) — an anniversary}

\label{lesson:TE-C355-varsikotsavam}



\begin{quote}
Before the new one: say the Telugu for to cover up, then the Telugu for a spy.
\end{quote}

\subsection*{You'll want to know: \te{వార్షికోత్సవం}}
\textbf{\te{వార్షికోత్సవం}} — \emph{vārṣikōtsavaṁ} — \textquotedblleft{}an anniversary\textquotedblright{}.

\begin{quote}
\te{మా} \te{బడి} \te{వార్షికోత్సవం} \te{ఘనంగా} \te{జరిగింది}.
\end{quote}

\begin{quote}
\emph{mā baḍi vārṣikōtsavaṁ ghanaṅgā jarigindi.}
\end{quote}

\begin{quote}
Our school's anniversary was celebrated in style.
\end{quote}

\te{వార్షిక}, yearly: a yearly celebration.

\begin{cousinweb}
\textbf{\te{వార్షికోత్సవం}} (\emph{vārṣikōtsavaṁ}) is borrowed from Sanskrit: \emph{vārṣika}, \textquotedblleft{}annual\textquotedblright{} (CDIAL 11571), and \emph{utsava}, \textquotedblleft{}a festival\textquotedblright{} (CDIAL 1876).
\end{cousinweb}

\begin{grammarlens}[title={What you've built}]
One more word for telling what happened.
\end{grammarlens}

\subsection*{Your turn}
\begin{itemize}
\raggedright
  \item \emph{Say it:} \emph{vārṣikōtsavaṁ}
  \item \emph{Say it:} \emph{vārṣikōtsavaṁ}, once more
  \item \emph{Say it:} \emph{mā baḍi vārṣikōtsavaṁ ghanaṅgā jarigindi.} — Our school's anniversary was celebrated in style.
  \item \emph{Recall:} say the Telugu for to cover up, then the Telugu for a spy, then say \emph{vārṣikōtsavaṁ} again
\end{itemize}

\begin{tcolorbox}[breakable,colback=teal!4,colframe=teal!35!black,title={Before you move on}]
What does \te{వార్షికోత్సవం} mean? (\textquotedblleft{}An anniversary\textquotedblright{}.) Say it once more.
\end{tcolorbox}
\section[\texorpdfstring{\te{పెళ్ళిరోజు} (peḷḷirōju)}{పెళ్ళిరోజు (peḷḷirōju)}]{\te{పెళ్ళిరోజు} (peḷḷirōju) — a wedding anniversary}

\label{lesson:TE-C355-pelliroju}



\begin{quote}
Before the new one: say the Telugu for a spy, then the Telugu for an anniversary.
\end{quote}

\subsection*{You'll want to know: \te{పెళ్ళిరోజు}}
\textbf{\te{పెళ్ళిరోజు}} — \emph{peḷḷirōju} — \textquotedblleft{}a wedding anniversary\textquotedblright{}.

\begin{quote}
\te{నిన్న} \te{అమ్మానాన్నల} \te{పెళ్ళిరోజు}.
\end{quote}

\begin{quote}
\emph{ninna ammānānnala peḷḷirōju.}
\end{quote}

\begin{quote}
Yesterday was my parents' wedding anniversary.
\end{quote}

\te{పెళ్ళి}, a wedding, and \te{రోజు}, a day: the day of the wedding, each year.

\begin{cousinweb}
\textbf{\te{పెళ్ళిరోజు}} (\emph{peḷḷirōju}) is a Telugu compound: native \emph{peḷḷi}, \textquotedblleft{}a marriage\textquotedblright{} (DEDR 4395), and \emph{rōju}, \textquotedblleft{}a day\textquotedblright{}.
\end{cousinweb}

\begin{grammarlens}[title={What you've built}]
One more word for telling what happened.
\end{grammarlens}

\subsection*{Your turn}
\begin{itemize}
\raggedright
  \item \emph{Say it:} \emph{peḷḷirōju}
  \item \emph{Say it:} \emph{peḷḷirōju}, once more
  \item \emph{Say it:} \emph{ninna ammānānnala peḷḷirōju.} — Yesterday was my parents' wedding anniversary.
  \item \emph{Recall:} say the Telugu for a spy, then the Telugu for an anniversary, then say \emph{peḷḷirōju} again
\end{itemize}

\begin{tcolorbox}[breakable,colback=teal!4,colframe=teal!35!black,title={Before you move on}]
What does \te{పెళ్ళిరోజు} mean? (\textquotedblleft{}A wedding anniversary\textquotedblright{}.) Say it once more.
\end{tcolorbox}
\section[\texorpdfstring{\te{షష్టిపూర్తి} (ṣaṣṭipūrti)}{షష్టిపూర్తి (ṣaṣṭipūrti)}]{\te{షష్టిపూర్తి} (ṣaṣṭipūrti) — a sixtieth birthday}

\label{lesson:TE-C355-sastipurti}



\begin{quote}
Before the new one: say the Telugu for an anniversary, then the Telugu for a wedding anniversary.
\end{quote}

\subsection*{You'll want to know: \te{షష్టిపూర్తి}}
\textbf{\te{షష్టిపూర్తి}} — \emph{ṣaṣṭipūrti} — \textquotedblleft{}a sixtieth birthday\textquotedblright{}.

\begin{quote}
\te{తాతగారి} \te{షష్టిపూర్తికి} \te{బంధువులు} \te{అందరూ} \te{వచ్చారు}.
\end{quote}

\begin{quote}
\emph{tātagāri ṣaṣṭipūrtiki bandhuvulu andarū vaccāru.}
\end{quote}

\begin{quote}
All the relatives came for grandfather's sixtieth birthday.
\end{quote}

\te{షష్టి}, sixty, and \te{పూర్తి}, completion: sixty years completed.

\begin{cousinweb}
\textbf{\te{షష్టిపూర్తి}} (\emph{ṣaṣṭipūrti}) is borrowed from Sanskrit: \emph{ṣaṣṭi}, \textquotedblleft{}sixty\textquotedblright{} (CDIAL 12804), and \emph{pūrti}, \textquotedblleft{}filling, completion\textquotedblright{} (Monier-Williams).
\end{cousinweb}

\begin{grammarlens}[title={What you've built}]
One more word for telling what happened.
\end{grammarlens}

\subsection*{Your turn}
\begin{itemize}
\raggedright
  \item \emph{Say it:} \emph{ṣaṣṭipūrti}
  \item \emph{Say it:} \emph{ṣaṣṭipūrti}, once more
  \item \emph{Say it:} \emph{tātagāri ṣaṣṭipūrtiki bandhuvulu andarū vaccāru.} — All the relatives came for grandfather's sixtieth birthday.
  \item \emph{Recall:} say the Telugu for an anniversary, then the Telugu for a wedding anniversary, then say \emph{ṣaṣṭipūrti} again
\end{itemize}

\begin{tcolorbox}[breakable,colback=teal!4,colframe=teal!35!black,title={Before you move on}]
What does \te{షష్టిపూర్తి} mean? (\textquotedblleft{}A sixtieth birthday\textquotedblright{}.) Say it once more.
\end{tcolorbox}
\section[\texorpdfstring{\te{జయంతి} (jayanti)}{జయంతి (jayanti)}]{\te{జయంతి} (jayanti) — a birth anniversary}

\label{lesson:TE-C355-jayanti}



\begin{quote}
Before the new one: say the Telugu for a wedding anniversary, then the Telugu for a sixtieth birthday.
\end{quote}

\subsection*{You'll want to know: \te{జయంతి}}
\textbf{\te{జయంతి}} — \emph{jayanti} — \textquotedblleft{}a birth anniversary\textquotedblright{}.

\begin{quote}
\te{ఆ} \te{కవి} \te{జయంతి} \te{సందర్భంగా} \te{సభ} \te{జరిగింది}.
\end{quote}

\begin{quote}
\emph{ā kavi jayanti sandarbhaṅgā sabha jarigindi.}
\end{quote}

\begin{quote}
A meeting was held for the poet's birth anniversary.
\end{quote}

For the famous: a poet's or a leader's \te{జయంతి}.

\begin{cousinweb}
\textbf{\te{జయంతి}} (\emph{jayanti}) is borrowed from Sanskrit \emph{jayantī} (Monier-Williams), a word built on \emph{jaya}, \textquotedblleft{}victory\textquotedblright{} (CDIAL 5141).
\end{cousinweb}

\begin{grammarlens}[title={What you've built}]
One more word for telling what happened.
\end{grammarlens}

\subsection*{Your turn}
\begin{itemize}
\raggedright
  \item \emph{Say it:} \emph{jayanti}
  \item \emph{Say it:} \emph{jayanti}, once more
  \item \emph{Say it:} \emph{ā kavi jayanti sandarbhaṅgā sabha jarigindi.} — A meeting was held for the poet's birth anniversary.
  \item \emph{Recall:} say the Telugu for a wedding anniversary, then the Telugu for a sixtieth birthday, then say \emph{jayanti} again
\end{itemize}

\begin{tcolorbox}[breakable,colback=teal!4,colframe=teal!35!black,title={Before you move on}]
What does \te{జయంతి} mean? (\textquotedblleft{}A birth anniversary\textquotedblright{}.) Say it once more.
\end{tcolorbox}
\section[\texorpdfstring{\te{సన్మానం} (sanmānaṁ)}{సన్మానం (sanmānaṁ)}]{\te{సన్మానం} (sanmānaṁ) — a felicitation; an honour}

\label{lesson:TE-C355-sanmanam}



\begin{quote}
Before the new one: say the Telugu for a sixtieth birthday, then the Telugu for a birth anniversary.
\end{quote}

\subsection*{You'll want to know: \te{సన్మానం}}
\textbf{\te{సన్మానం}} — \emph{sanmānaṁ} — \textquotedblleft{}a felicitation; an honour\textquotedblright{}.

\begin{quote}
\te{సభలో} \te{పాత} \te{మాస్టారికి} \te{సన్మానం} \te{జరిగింది}.
\end{quote}

\begin{quote}
\emph{sabhalō pāta māsṭāriki sanmānaṁ jarigindi.}
\end{quote}

\begin{quote}
The old teacher was felicitated at the meeting.
\end{quote}

\te{సన్మానం} \te{చేయు}, to felicitate: a shawl, a garland and a speech.

\begin{cousinweb}
\textbf{\te{సన్మానం}} (\emph{sanmānaṁ}) is borrowed from Sanskrit \emph{sanmāna}, \textquotedblleft{}respect, esteem\textquotedblright{} (Monier-Williams).
\end{cousinweb}

\begin{grammarlens}[title={What you've built}]
One more word for telling what happened.
\end{grammarlens}

\subsection*{Your turn}
\begin{itemize}
\raggedright
  \item \emph{Say it:} \emph{sanmānaṁ}
  \item \emph{Say it:} \emph{sanmānaṁ}, once more
  \item \emph{Say it:} \emph{sabhalō pāta māsṭāriki sanmānaṁ jarigindi.} — The old teacher was felicitated at the meeting.
  \item \emph{Recall:} say the Telugu for a sixtieth birthday, then the Telugu for a birth anniversary, then say \emph{sanmānaṁ} again
\end{itemize}

\begin{tcolorbox}[breakable,colback=teal!4,colframe=teal!35!black,title={Before you move on}]
What does \te{సన్మానం} mean? (\textquotedblleft{}A felicitation; an honour\textquotedblright{}.) Say it once more.
\end{tcolorbox}
